\documentclass[10pt,a4paper]{amsart}
\usepackage[margin=19mm]{geometry}
\usepackage{amsmath,amssymb,mathtools}
\usepackage[scr=rsfs,cal=euler]{mathalfa}
\usepackage{xfrac}
\usepackage[unicode,hidelinks,bookmarks=false]{hyperref}
\newcommand{\ie}{i.e.\ }
\newcommand{\cf}{cf.\ }
\newcommand{\resp}{respectively\ }
\newcommand{\Subsection}{\subsection}
\providecommand{\nobreakdash}{\nobreak\hbox{-}}
\setlength{\emergencystretch}{2em}
\multlinegap0pt
\usepackage{xcolor}
\usepackage{amssymb}
\usepackage{tikz}
\newtheorem{prop}{Proposition}
\newcommand{\La}{\Lambda}
\newcommand{\Pol}{\mathrm{Pol}}
\newcommand{\ddd}{\mathfrak{d}}
\newcommand{\kb}{\kappa}
\newcommand{\nd}{n_\bullet}
\newcommand{\zd}{z_\bullet}
\title{Wreath Macdonald operators}
\author{Daniel Orr, Mark Shimozono, Joshua Jeishing Wen}
\date{Source: arXiv:2211.03851v3 (CC BY 4.0)}
\begin{document}
\maketitle

\noindent\emph{Source and licence.} Wreath Macdonald operators. Version arXiv:2211.03851v3 (CC BY 4.0), distributed by arXiv under the Creative Commons Attribution 4.0 International licence, http://creativecommons.org/licenses/by/4.0/, as recorded in the arXiv licence metadata for this exact version and shown on the abstract page https://arxiv.org/abs/2211.03851v3. Full licence notice: \texttt{licenses/arxiv-2211.03851-CC-BY-4.0.txt}.\par\smallskip

\noindent\textit{Editorial scope.} Counted unit: the article's prop FiniteVert (source0.tex lines 2149--2190). The article's complete original proof of it is delivered with it (source0.tex lines 2192--2263, one \texttt{proof} environment of 72 lines). No mathematics is written, repaired or re-derived: the statement and the proof are the article's own lines, in the article's own notation, and no retained line is substituted or re-flowed.  Retained uncounted as the source-local context the proof needs: lines 947--949; lines 1359--1414.  Nothing was omitted from any retained span, so the package carries no omission record.  No retained line contains a citation, so the wrapper carries no bibliography and no bibliography entry.  No retained line contains a figure environment, an image-inclusion command or drawing code, so no image is delivered and the package declares no asset dependency.  Editorial formatting only, in addition: an AMS \texttt{amsart} preamble that restates the article's own declarations and macros used by the retained lines, namely the environments \texttt{prop} on separate counters, together with the standard packages those lines need.  The retained environments print in this document as Proposition 1 in order of appearance, whereas the article prints the same items with the numbering of its own full document.  The complete native source of the article as distributed in the arXiv e-print for this exact version is bundled unchanged as \texttt{sources/133143/source0.tex}.
\begin{align}\label{Compatibility}
N_{i}-N_{i-1}=(\alpha_i^\vee,\kb(\lambda))=(\alpha_i^\vee,\alpha)=c_{i-1}-c_{i},\quad \text{for all $i\in I$,}
\end{align}

\subsubsection{Normal ordering}
Later, we will make use of a particular expression for products of the currents $\{E_i(z)\}$ and $\{F_i(z)\}$.
We will need notation for an ordered product or composition of noncommuting operators $a_1,\ldots, a_m$:
\begin{equation}
    \begin{aligned}
    \overset{\curvearrowright}{\prod_{j=1}^m}a_j&:=a_1 a_2\cdots a_m,&
    \overset{\curvearrowleft}{\prod_{j=1}^m}a_j&:=a_m a_{m-1}\cdots a_1
    \end{aligned}
\label{OrdProd}
\end{equation}
\begin{prop}\label{NormalOrder}
For $p\in I$, we have
\begin{equation}
\begin{aligned}
&\overset{\curvearrowright}{\prod_{a=1}^n}\,\overset{\curvearrowright}{\prod_{i=1}^{r}}E_{p+i}(z_{p+i,a})\\
&=
\left( (-1)^{\frac{(r-2)(r-3)}{2}}\ddd^{\frac{r}{2}-1}\prod_{i\in I} c_i\right)^n\\
&\times\prod_{1\le a<b\le n}\prod_{i\in I}
\frac{\displaystyle\left(1-z_{i,b}/z_{i,a}\right)\left(1-q^{-1}t^{-1}z_{i,b}/z_{i,a}\right)}{\displaystyle\left(1-t^{-1}z_{i+1,b}/z_{i,a}\right)\left(1-q^{-1}z_{i-1,b}/z_{i,a}\right)}\\
&\times\prod_{a=1}^n \frac{z_{p,a}/z_{p+1,a}}
{\left(1-q^{-1}z_{p,a}/z_{p+1,a}\right)
\prod_{i\in I\setminus\{p+1\}}
\left(1-t^{-1}z_{i,a}/z_{i-1,a}\right)}\\
&\times\prod_{i\in I}\exp\left(\sum_{a=1}^n\sum_{k>0}\left(p_{k}[X^{(i)}]-t^{-k}p_k[X^{(i-1)}]\right)\frac{z_{i,a}^{k}}{k}\right)\\
&\times\prod_{i\in I}\exp\left(\sum_{a=1}^n\sum_{k>0}\left(-p_{k}[X^{(i)}]^\perp+q^{-k}p_k[X^{(i-1)}]^\perp\right)\frac{z_{i,a}^{-k}}{k}\right)
\prod_{i\in I}\prod_{a=1}^n z_{i,a}^{H_{i,0}}
\end{aligned}
\label{ENormalOrder}
\end{equation}
where all rational functions are Laurent series expanded assuming 
\begin{equation}
|z_{i,a}|=1,\,|q|>1,\,|t|>1.
\label{NormalEExpand}
\end{equation}
For the $F$-currents, we have
\begin{equation}
\begin{aligned}
&\overset{\curvearrowleft}{\prod_{a=1}^n}\,\overset{\curvearrowleft}{\prod_{i=1}^{r}}F_{p+i}(z_{p+i,a})\\
&=\left( \frac{(-1)^{\frac{(r-2)(r-3)}{2}}}{\displaystyle\ddd^{\frac{r}{2}-1} \prod_{i\in I} c_i}\right)^n\\
&\times\prod_{1\le a<b\le n}\prod_{i\in I}
\frac{\left(1-z_{i,a}/z_{i,b}\right)\left(1-qtz_{i,a}/z_{i,b}\right)}{\left(1-tz_{i-1,a}/z_{i,b}\right)\left(1-qz_{i+1,a}/z_{i,b}\right)}\\
&\times \prod_{a=1}^n\frac{z_{p+1,a}/z_{p,a}}{\left(1-qz_{p+1,a}/z_{p,a}\right)
\prod_{i\in I\setminus\{p\}}
\left(1-tz_{i,a}/z_{i+1,a}\right)}\\
&\times\prod_{i\in I}\exp\left(\sum_{a=1}^n\sum_{k>0}\left(-t^{k}p_{k}[X^{(i)}]+p_k[X^{(i-1)}]\right)\frac{z_{i,a}^{k}}{k}\right)\\
&\times\prod_{i\in I}\exp\left(\sum_{a=1}^n\sum_{k>0}\left(q^{k}p_{k}[X^{(i)}]^\perp-p_k[X^{(i-1)}]^\perp\right)\frac{z_{i,a}^{-k}}{k}\right)
\prod_{i\in I}\prod_{a=1}^n z_{i,a}^{-H_{i,0}}
\end{aligned}
\label{FNormalOrder}
\end{equation}
where all rational functions are Laurent series expanded assuming 
\begin{equation}
|z_{i,a}|=1,\, |q|< 1,\, |t|<1.
\label{NormalFExpand}
\end{equation}
\end{prop}

\begin{prop}\label{FiniteVert}
    Let $f\in\Lambda^I$ be factored according to color:
    \[
    f= \prod_{i\in I} f_i[X^{(i)}],
    \]
    where $f_i\in\Lambda$ for all $i\in I$.
    For
    \begin{equation}
    |z|=1,\, |q|>1,\, |t|>1,\, |x_{l}^{(j)}|<1,
    \label{EExpand0}
    \end{equation}
    the vertex operators from (\ref{ENormalOrder}) act on $f$ such that upon finitization, we have:
\begin{equation}
\label{FinEVertex}
    \begin{aligned}
    &\pi_{N_\bullet} \left(\exp\left[ \sum_{k>0}\left( p_k[X^{(i)}]-t^{-k}p_k[X^{(i-1)}] \right) \frac{z^k}{k}\right]\right.\\
    &\times\left.\exp\left[\sum_{k>0}\left( -p_k[X^{(i)}]^\perp+q^{-k}p_k[X^{(i-1)}]^\perp \right)\frac{z^{-k}}{k}  \right]z^{H_{i,0}} (f\otimes e^\alpha)\right) \\
    &=
    \frac{\displaystyle\prod_{l=1}^{N_{i-1}}\left(z^{-1}-t^{-1}x_{l}^{(i-1)}\right)}{\displaystyle\prod_{l=1}^{N_i} \left(z^{-1}-x_{l}^{(i)}\right)}
    \left( f_i\left[x_{\bullet}^{(i)}-z^{-1}\right]f_{i-1}\left[x_{\bullet}^{(i-1)}+q^{-1}z^{-1}\right]
    \prod_{\substack{j\in I\\j\not=i,i-1}} f_j\left[x_{\bullet}^{(j)}\right]\otimes e^\alpha\right).
    \end{aligned}
\end{equation}
    On the other hand, for
    \begin{equation*}
    |z|=1,\, |q|<1,\, |t|<1,\, |x_{l}^{(j)}|<1,
    \label{FExpand0}
    \end{equation*}
    the vertex operators from (\ref{FNormalOrder}) act as:
\begin{equation}
\label{FinFVertex}
    \begin{aligned}
    &\pi_{N_\bullet} \left(\exp\left[ \sum_{k>0}\left( -t^kp_k[X^{(i)}]+p_k[X^{(i-1)}] \right) \frac{z^k}{k}\right]\right.\\
    &\times \left.\exp\left[\sum_{k>0}\left( q^kp_k[X^{(i)}]^\perp-p_k[X^{(i-1)}]^\perp \right)\frac{z^{-k}}{k}  \right]z^{-H_{i,0}} (f\otimes e^\alpha)\right) \\
    &=
    \frac{\displaystyle\prod_{l=1}^{N_{i}}\left(z^{-1}-tx_{l}^{(i)}\right)}
    {\displaystyle\prod_{l=1}^{N_{i-1}} \left(z^{-1}-x_{l}^{(i-1)}\right)}
    \left( f_i\left[x_{\bullet}^{(i)}+qz^{-1}\right]f_{i-1}\left[x_{\bullet}^{(i-1)}-z^{-1}\right]
    \prod_{\substack{j\in I\\j\not=i,i-1}} f_j\left[x_{\bullet}^{(j)}\right]\otimes e^\alpha\right).
    \end{aligned}
\end{equation}
\end{prop}

\begin{proof}
We will only consider (\ref{FinEVertex})---the proof for (\ref{FinFVertex}) is similar.
First consider the `left' half of the vertex operator together with $z^{H_i,0}$. 
We have
\begin{align}
\nonumber
&\pi_{N_\bullet} \left(\exp\left[ \sum_{k>0}\left( p_k[X^{(i)}]-t^{-k}p_k[X^{(i-1)}] \right) \frac{z^k}{k}\right]z^{H_{i,0}} (f\otimes e^\alpha)\right) \\
\label{FinCreation}
&=
\frac{\displaystyle\prod_{l=1}^{N_{i-1}}\left(1-t^{-1}zx_{l}^{(i-1)}\right)}{\displaystyle\prod_{l=1}^{N_i} \left(1-zx_{l}^{(i)}\right)}
z^{N_i-N_{i-1}}
\left( \pi_{N_\bullet}(f)\otimes e^\alpha\right)
=
\frac{\displaystyle\prod_{l=1}^{N_{i-1}}\left(z^{-1}-t^{-1}x_{l}^{(i-1)}\right)}{\displaystyle\prod_{l=1}^{N_i} \left(z^{-1}-x_{l}^{(i)}\right)}
\left( \pi_{N_\bullet}(f)\otimes e^\alpha\right)
\end{align}
%for any $f\in\Lambda^{\otimes I}$.
For this to hold, we will need to impose conditions on $|x_{l}^{(i)}|$.
Recall that we have the conditions (\ref{NormalEExpand}) when working with $\{E_i(z)\}$.
We extend these to (\ref{EExpand0}) for (\ref{FinCreation}) to hold. 
Let us also point out that the compatibility condition \eqref{Compatibility} is used to obtain the factor $z^{N_i-N_{i-1}}$ after the first equality.

Next, from the `right' half, we have
\begin{align}
\nonumber
&\pi_{N_\bullet}\left( \exp\left[\sum_{k>0}\left( -p_k[X^{(i)}]^\perp+q^{-k}p_k[X^{(i-1)}]^\perp \right)\frac{z^{-k}}{k}  \right]\cdot f \right)\\
\label{FinAnnihilation}
&=
\pi_{N_\bullet}\left(f_i[X^{(i)}-z^{-1}]f_{i-1}[X^{(i-1)}+q^{-1}z^{-1}]\prod_{\substack{j\in I\\j\not=i,i-1}} f_j[X^{(j)}] \right)\\
\nonumber
&= f_i\left[x_{\bullet}^{(i)}-z^{-1}\right]f_{i-1}\left[x_{\bullet}^{(i-1)}+q^{-1}z^{-1}\right]
\prod_{\substack{j\in I\\j\not=i,i-1}} f_j\left[x_{\bullet}^{(j)}\right].\notag
\end{align}
Here, (\ref{FinAnnihilation}) follows from checking on power sums $p_k[X^{(i)}]$ and $p_k[X^{(i-1)}]$.
%\begin{align*}
%\pi_n\big(p_m[X^{(i)}-z_i^{-1}+q^{-1}z_{i+1}^{-1}]\big)
%&= \pi_n\big(p_m[X^{(i)}]-z_i^{-m}+q^{-m}z_{i+1}^{-m}]\big)\\
%&= p_m[x^{(i)}_1+\dotsm +x^{(i)}_{n_i}]-z_i^{-m}+q^{-m}z_{i+1}^{-m}\\
%&= p_m[x^{(i)}_1+\dotsm +x^{(i)}_{n_i}-z_i^{-1}+q^{-1}z_{i+1}^{-1}].
%\end{align*}
%Let us denote by $\pi_{\nd}^{\zd} : \La^I \to \Pol_{\nd}[z_0^{-1},z_1^{-1},\dotsc,z_{r-1}^{-1}]$ the map given by the simultaneous plethystic substitutions $X^{(i)}\mapsto x^{(i)}_1+\dotsm +x^{(i)}_{n_i}-z_i^{-1}+q^{-1}z_{i+1}^{-1}$ for each $i\in I$, i.e.,
%\begin{align}
%\pi_{\nd}^{\zd}(f) &= \pi_{\nd}\Big(\prod_{j\in I}\Omega[-z_j^{-1}(X^{(j)}-q^{-1}X^{(j-1)})]^\perp\cdot f  \Big).
%\end{align}
%
%For \eqref{FNormalOrder}, we have
%\begin{align*}
%\nonumber
%&\pi_{N_\bullet} \left(\exp\left[ \sum_{k>0}\left( -t^kp_k[X^{(i)}]+p_k[X^{(i-1)}] \right) \frac{z^k}{k}\right]z^{-H_{i,0}} (f\otimes e^\alpha)\right) \\
%&=
%\frac{\displaystyle\prod_{l=1}^{N_{i}}\left(1-tzx_{l}^{(i)}\right)}
%{\displaystyle\prod_{l=1}^{N_{i-1}} \left(1-zx_{l}^{(i-1)}\right)}
%z^{N_{i-1}-N_{i}}
%\left( \pi_{N_\bullet}(f)\otimes e^\alpha\right)
%=
%\frac{\displaystyle\prod_{l=1}^{N_{i}}\left(z^{-1}-tx_{l}^{(i)}\right)}
%{\displaystyle\prod_{l=1}^{N_{i-1}} \left(z^{-1}-x_{l}^{(i-1)}\right)}
%\left( \pi_{N_\bullet}(f)\otimes e^\alpha\right).
%\end{align*}
%Here, we extend (\ref{NormalFExpand}) to
%
%On the other hand, from the right halves appearing in \eqref{FNormalOrder}, we have
%\begin{align*}
%\nonumber
%&\pi_{N_\bullet}\left( \exp\left[\sum_{k>0}\left( q^kp_k[X^{(i)}]^\perp-p_k[X^{(i-1)}]^\perp \right)\frac{z^{-k}}{k}  \right]\cdot f  \right)\\
%&=
%\pi_{N_\bullet}\left(f_i[X^{(i)}+qz^{-1}]f_{i-1}[X^{(i-1)}-z^{-1}]\prod_{\substack{j\in I\\j\not=i,i-1}} f_j[X^{(j)}]  \right)\\
%\nonumber
%&= f_i\left[x_{\bullet}^{(i)}+qz^{-1}\right]f_{i-1}\left[x_{\bullet}^{(i-1)}-z^{-1}\right]
%\prod_{\substack{j\in I\\j\not=i,i-1}} f_j\left[x_{\bullet}^{(j)}\right].\notag
%\end{align*}
\end{proof}

\end{document}
